\section{Reduction rules}
\label{app:reduction}

This appendix gives the complete reduction rules of the Statix type checker
summarised in \cref{sec:beta-reduction}.

\subsection{Weak head normalisation}

Weak head normalisation is implemented as a head-reduction procedure with an
explicit argument spine. We write
\[
\mathsf{whnf}_h(s; e; \overline{a}) = (s', e')
\]
to mean that expression \(e\), evaluated in scope \(s\), reduces to weak head
normal form \(e'\) under scope \(s'\), while carrying a pending argument spine
\(\overline{a}\). Each argument in the spine is a closure \((s_i,e_i)\), where
\(s_i\) is the scope in which the argument \(e_i\) must be interpreted. The
ordinary weak head normalisation function is the special case with an empty
argument spine, \(\mathsf{whnf}(s;e) = \mathsf{whnf}_h(s;e;\epsilon)\).
\Cref{fig:whnf} gives the rules.

\begin{figure}[tbp]
\[
\begin{array}{rcll}
\mathsf{whnf}_h(s; e_1\ e_2; \overline{a})
&=&
\mathsf{whnf}_h(s; e_1; (s,e_2) :: \overline{a})
& \textsc{Wh-App}
\\[0.8em]

\mathsf{whnf}_h(s; (e : A); \overline{a})
&=&
\mathsf{whnf}_h(s; e; \overline{a})
& \textsc{Wh-Ann}
\\[0.8em]

\mathsf{whnf}_h(s; \mathbf{let}\ x = v\ \mathbf{in}\ b; \overline{a})
&=&
\mathsf{whnf}_h(s[x' \mapsto (s,v)]; b; \overline{a})
& \textsc{Wh-Let}
\\
&&
\multicolumn{2}{l}{\text{where } x' = \mathsf{fresh}(s,x)}
\\[0.8em]

\mathsf{whnf}_h(s; \lambda x.\ b; (s_a,a) :: \overline{a})
&=&
\mathsf{whnf}_h(s[x' \mapsto (s_a,a)]; b; \overline{a})
& \textsc{Wh-Beta}
\\
&&
\multicolumn{2}{l}{\text{where } x' = \mathsf{fresh}(s,x)}
\\[0.8em]

\mathsf{whnf}_h(s; \lambda x.\ b; \epsilon)
&=&
(s,\lambda x.\ b)
& \textsc{Wh-Lam}
\\[0.8em]

\mathsf{whnf}_h(s; H; \overline{a})
&=&
\mathsf{rebuild}((s,H), \overline{a})
& \textsc{Wh-Head}
\\
&&
\multicolumn{2}{l}{
H \in
\{
\Pi x : A.\ B,\ A \to B,\ \mathsf{Set}
\}}
\\[0.8em]

\mathsf{whnf}_h(s; x; \overline{a})
&=&
\mathsf{whnf}_h(s_v; v; \overline{a})
& \textsc{Wh-VarSubst}
\\
&&
\multicolumn{2}{l}{
\text{if } \mathsf{lookupSubst}(s,x) = (s_v,v)}
\\[0.8em]

\mathsf{whnf}_h(s; x; \overline{a})
&=&
\mathsf{rebuild}((s,x), \overline{a})
& \textsc{Wh-VarStuck}
\\
&&
\multicolumn{2}{l}{
\text{if } \mathsf{lookupSubst}(s,x) \text{ is undefined}}
\end{array}
\]
\caption{Weak head normalisation with an explicit argument spine.}
\label{fig:whnf}
\Description{Eight equations defining whnf_h: Wh-App, Wh-Ann, Wh-Let and Wh-Beta perform head computation, Wh-Lam leaves a lambda without pending arguments unchanged, and Wh-Head, Wh-VarSubst and Wh-VarStuck handle stuck heads and variables with or without a delayed substitution.}
\end{figure}

The first four rules perform head computation. Applications are decomposed into
a head and an explicit argument spine. Type annotations are ignored during
reduction. Let-bindings extend the current scope with a delayed substitution,
and lambda abstractions consume one pending argument by extending the current
scope with a delayed substitution for the bound variable. If a lambda has no
pending argument, it is already in weak head normal form.

The remaining rules describe stuck heads. A \(\Pi\)-type and \(\mathsf{Set}\)
do not reduce at the head position, so the pending spine is rebuilt. Variables
reduce only when they have a delayed substitution in the current scope.
Otherwise, they are neutral terms, and the spine is rebuilt around them.
\(\mathsf{rebuild}\) reconstructs applications once head reduction becomes
stuck:
\[
\begin{array}{rcl}
\mathsf{rebuild}((s,e), \epsilon)
&=&
(s,e)
\\[0.5em]
\mathsf{rebuild}((s,e), a :: \overline{a})
&=&
\mathsf{rebuild}(\mathsf{apply}((s,e),a),\overline{a}).
\end{array}
\]
Applying a head to an argument preserves the argument's original scope. If the
head and argument are already in the same scope, the application can be rebuilt
directly:
\[
\mathsf{apply}((s,e),(s,a)) = (s, \mathsf{App}(e, a)).
\]
If the argument comes from a different scope or we have made some progress
evaluating it, the argument is stored as a fresh temporary substitution, and the
rebuilt application refers to the temporary variable:
\[
\mathsf{apply}((s,e),(s_a,a))
=
(s[t \mapsto (s_a,a)],\ \mathsf{App}(e, t))
\qquad
\text{where } t = \mathsf{fresh}(\_, \_app).
\]

\subsection{Full normalisation}

Full normalisation first reduces the input expression to weak head normal form
and then, once the outermost constructor is known, recursively normalises its
immediate subexpressions. Unlike weak head normalisation, it continues under
lambda abstractions, inside applications, and inside function types. We write
\(\mathsf{normalise}(s,e) = e'\) to mean that expression \(e\), evaluated in
scope \(s\), fully normalises to \(e'\). The auxiliary judgment
\(\mathsf{normalise}_{\mathsf{wh}}(s,e) = e'\) normalises an expression that has
already been reduced to weak head normal form. \Cref{fig:normalisation} gives
the rules.

\begin{figure}[tbp]
  \begin{mathpar}
  \inferrule*[right=\textsc{Norm-Whnf}]
    {
      \mathsf{whnf}(s,e) = (s',w)
      \\
      \mathsf{normalise}_{\mathsf{wh}}(s',w) = n
    }
    {
      \mathsf{normalise}(s,e) = n
    }

  \and

  \inferrule*[right=\textsc{Norm-Pi}]
    {
      \mathsf{normalise}(s,A) = A'
      \\
      \mathsf{normalise}(s,B) = B'
    }
    {
      \mathsf{normalise}_{\mathsf{wh}}(s,\Pi x : A.\ B)
      =
      \Pi x : A'.\ B'
    }

  \and

  \inferrule*[right=\textsc{Norm-Lam}]
    {
      \mathsf{normalise}(s,b) = b'
    }
    {
      \mathsf{normalise}_{\mathsf{wh}}(s,\lambda x.\ b)
      =
      \lambda x.\ b'
    }

  \and

  \inferrule*[right=\textsc{Norm-App}]
    {
      \mathsf{normalise}(s,f) = f'
      \\
      \mathsf{normalise}(s,a) = a'
    }
    {
      \mathsf{normalise}_{\mathsf{wh}}(s,f\ a)
      =
      f'\ a'
    }

  \and

  \inferrule*[right=\textsc{Norm-Neutral}]
    { }
    {
      \mathsf{normalise}_{\mathsf{wh}}(s,e) = e
    }

  \end{mathpar}
  \caption{Full normalisation rules.}
  \label{fig:normalisation}
  \Description{Inference rules Norm-Whnf, Norm-Pi, Norm-Lam, Norm-App and Norm-Neutral, defining full normalisation of an expression in a scope via weak head normalisation followed by normalisation of the immediate subexpressions.}
\end{figure}
